% ST图原理示意图
\begin{tikzpicture}[scale=0.85, transform shape]
  % 坐标轴
  \draw[-Stealth, thick] (0,0) -- (9,0) node[right, font=\small\bfseries] {$t$ (s)};
  \draw[-Stealth, thick] (0,0) -- (0,7.5) node[above, font=\small\bfseries] {$s$ (m)};

  % 刻度
  \foreach \t in {1,...,7} {
    \draw (\t, -0.1) -- (\t, 0.1);
    \node[below, font=\footnotesize] at (\t, -0.1) {\t};
  }
  \foreach \s in {1,...,7} {
    \draw (-0.1, \s) -- (0.1, \s);
    \node[left, font=\footnotesize] at (-0.1, \s) {\pgfmathparse{int(\s*5)}\pgfmathresult};
  }

  % 静态障碍物ST边界（矩形）
  \fill[dangerred!30, draw=dangerred, thick] (0,2.5) rectangle (7.5,3.5);
  \node[font=\footnotesize\bfseries, dangerred] at (3.8,3.0) {静态障碍物};

  % 动态障碍物ST边界（倾斜平行四边形）
  \fill[lightorange!60, draw=lightorange!80!black, thick]
    (1,5.5) -- (5,6.5) -- (5,7.2) -- (1,6.2) -- cycle;
  \node[font=\footnotesize\bfseries, lightorange!80!black] at (3,6.6) {动态障碍物};

  % 可行轨迹（绿色曲线 - 从障碍物下方通过）
  \draw[very thick, lightgreen!70!black, -{Stealth}]
    plot[smooth, tension=0.5] coordinates {(0,0.5) (1,1.2) (2,1.8) (3,2.1) (4,2.2) (5,2.3) (6,2.3) (7,2.4)};
  \node[lightgreen!70!black, font=\footnotesize\bfseries] at (7.5,2.7) {轨迹1(YIELD)};

  % 可行轨迹2（蓝色曲线 - 从障碍物上方通过）
  \draw[very thick, deepblue, dashed, -{Stealth}]
    plot[smooth, tension=0.5] coordinates {(0,0.5) (0.8,1.8) (1.5,3.0) (2,4.0) (3,4.8) (4,5.3) (5,5.5) (6,5.8) (7,6.0)};
  \node[deepblue, font=\footnotesize\bfseries] at (7.5,6.3) {轨迹2(OVERTAKE)};

  % 安全距离标注
  \draw[{Stealth}-{Stealth}, thick, deeppink] (7.8, 2.5) -- (7.8, 3.5);
  \node[deeppink, font=\tiny, right] at (7.9, 3.0) {$d_{ext}$};

  % 不可行区域标注
  \draw[-Stealth, thick, gray] (5.5, 4.5) -- (4.5, 3.3);
  \node[gray, font=\footnotesize, right] at (5.2, 4.7) {禁行区域};

  % 图例
  \fill[dangerred!30, draw=dangerred] (0.5,-1.0) rectangle (1.2,-0.6);
  \node[right, font=\footnotesize] at (1.3,-0.8) {静态障碍物ST边界};
  \fill[lightorange!60, draw=lightorange!80!black] (4.5,-1.0) rectangle (5.2,-0.6);
  \node[right, font=\footnotesize] at (5.3,-0.8) {动态障碍物ST边界};
\end{tikzpicture}
